\documentclass[11pt, oneside]{article} 
\usepackage{geometry}
\geometry{letterpaper} 
\usepackage{graphicx}
	
\usepackage{amssymb}
\usepackage{amsmath}
\usepackage{parskip}
\usepackage{color}
\usepackage{hyperref}

\graphicspath{{/Users/telliott/Dropbox/Github-math/figures/}}
% \begin{center} \includegraphics [scale=0.4] {gauss3.png} \end{center}

\title{Introduction}
\date{}

\begin{document}
\maketitle
\Large

%[my-super-duper-separator]

In this chapter we look at the Simson line of a given triangle:

\url{https://en.wikipedia.org/wiki/Simson_line}

Find an arbitrary point $P$ on the circumcircle of $\triangle ABC$.  $PCAB$ is a cyclic quadrilateral.

\begin{center} \includegraphics [scale=0.3] {simson1.png} \end{center}

Drop perpendiculars from $P$ to the sides of the triangle (or their extensions).  The points of intersection, $L,M,N$, are collinear.

\[ \angle PBA + \angle ACP = 180 \]
\[ \angle PBA = \angle PBN \]
\[ \angle PBN + \angle ACP = 180 \]

Because of the right angles, $PLAM$ and $PMNC$ are also cyclic.

\begin{center} \includegraphics [scale=0.35] {simson3.png} \end{center}

Here is the wikipedia proof.

\emph{Proof}.

From the circle with $PB$ as diameter (right panel):
\[ \angle PBN + \angle NMP = 180 \]
Combining with the first result above: $ \angle NMP = \angle ACP$.

In the circle with $PC$ as diameter (left panel), $\angle PML$ and $\angle PCL$ lie on the same arc so
\[ \angle PML = \angle PCL = 180 - \angle ACP \]

Finally, 
\[ \angle PML = 180 - \angle NMP \]
so $L,M,N$ are collinear.

$\square$

Here is my shortest proof.

\emph{Proof}.

The supplementary angle to $\angle ACP$ is $\angle LCP$ (left panel).

But another angle supplementary to $\angle ACP$ is $\angle ABP$ (original quadrilateral).

\begin{center} \includegraphics [scale=0.35] {simson3.png} \end{center}

$\angle LPC$ is complementary to the first, and $\angle BPN$ is complementary to the second, hence $\angle LPC = \angle BPN$.

But $\angle LMC$ and $\angle LPC$ lie on equal arcs of the circle on $PC$, while $\angle BMN$ and $\angle BPN$ lie equal arcs of the circle on $PB$.

It follows that $\angle BMN = \angle LMC$.

Therefore, $L,M,N$ are collinear.

$\square$

It seems noteworthy that the angles at $P$, $\angle LPM = C$ and $\angle NPM = B$, but I haven't been able to use that to get further.

Coxeter use the Simson line to develop a proof of Ptolemy's theorem.  It uses a preliminary result about \emph{pedal} triangles that we haven't seen before.

\subsection*{pedal triangles}

In the interior of $\triangle ABC$ find an arbitrary point $P$ and drop perpendiculars to the sides at $L,M,N$.

The pedal triangle is $\triangle LMN$.

\begin{center} \includegraphics [scale=0.35] {pedal.png} \end{center}

$PC$ is the diameter of a circle on which $M$ and $N$ lie.  By the law of sines, $MN/\sin C = PC$.

But in $\triangle ABC$, $c/\sin C = 2R$ for the same reason.

\[ \frac{MN}{PC} = \frac{c}{2R} \]
\[ MN = \frac{c \cdot PC}{2R} \]

and similarly for the other sides.  

\subsection*{Ptolemy's theorem}

Where it gets weird is that Coxeter says that in "triangle" $L,M,N$ this still holds, even though the triangle is degenerate.

\begin{center} \includegraphics [scale=0.3] {simson1.png} \end{center}

Remember, we dropped perpendiculars from $P$ to two sides of $\triangle ABC$ and joined those points.  So $P$ to $L$ on side $AC$ and $P$ to $M$ on side $BC$.

Then the length of the third side of $\triangle ABC$, $AB$, is related to the distance from $P$ across the pedal side to a vertex of the original triangle, opposite to the side in question.

That would be $PC$.
\[  LM = AB \cdot PC/2R \]

The other two are 

\[ LN = BC \cdot PA/2R \]
\[ MN = AC \cdot PB/2R \]

And since 
\[ LM + MN = LN \]

Finally
\[ AB \cdot PC + AC \cdot PB = BC \cdot PA  \]
which is Ptolemy's theorem.

I'm not totally convinced, but I'm going to add

$\square$

\end{document}